\documentclass[11pt,letterpaper,twoside]{article}
\usepackage[margin=.9in,headheight=40pt]{geometry}
\usepackage[kripke]{andyspackage}
\usepackage[header=book,font=times,accent=purple]{andysnotes}
\title{Andy's Package}
\AndySubtitle{Quick start and notation guide}
\author{Andrew M. Hagy}
\date{September 2026}
\begin{document}
\setcounter{page}{0}
\AndyMakeTitle

\section{Start a document}
Copy the master Overleaf project. Its root contains \texttt{main.tex}
and the three style files. Edit \texttt{main.tex}; there is no need to
import the packages into every new project.

Select one notation line and one style line in the preamble. A percent
sign comments out a line; move it so that exactly one style line is active.
The default is a portable Times-family face, purple accents, and
alternating running headers on facing pages. Overleaf's pdfLaTeX
compiler works for this default.
\begin{Verbatim}[fontsize=\small]
\usepackage{andyspackage}
% \usepackage[kripke]{andyspackage}
% \usepackage[kripke,prooftrees]{andyspackage}

\usepackage[header=book,font=times]{andysnotes}
% \usepackage[header=plain,font=times]{andysnotes}
% \usepackage[header=none,font=times]{andysnotes}
% \usepackage{andyspics}
\end{Verbatim}
The \texttt{book} header alternates a running title and the page
number with a double rule; \texttt{plain} uses a simple running header;
\texttt{none} removes both. The first content page is numbered 0
by the starter file. Change \verb|\setcounter{page}{0}| to 1 if
a conventional first page suits the project.

\section{Choose only the title fields you need}
A title and author suffice. The subtitle, course, assignment, due
date, and affiliation are independent optional commands. Use
\verb|\AndyMakeTitle| for a title above the first section, or
\verb|\AndyMakeTitlePage| for a separate cover. Uncomment
\verb|\tableofcontents| only when useful.
\begin{Verbatim}[fontsize=\small]
\title{A Working Title}
\author{Andrew M. Hagy}
\date{}
% \AndyCourse{PHIL 4770}
% \AndyDue{Friday, October 2, 2026}
\begin{document}
\setcounter{page}{0}
\AndyMakeTitle
% \tableofcontents
\section{Introduction}
\subsection{First idea}
\end{document}
\end{Verbatim}
Section headings are unnumbered on the page, but they advance the
section counter. Thus the first subsection prints \textbf{1.1},
and the first definition there prints \textbf{1.1.1}. A definition,
lemma, proposition, claim, example, remark, note, exercise, chaff,
or aside takes the next number in this shared sequence. Theorems and
conjectures have separate global sequences. A corollary after
Theorem 1 prints \textbf{Corollary 1.i}.

\section{Write with the notation}
\[
  \Gamma\proves\varphi,\qquad
  \Gamma\entails\varphi,\qquad
  \Str\satisfies\varphi,\qquad
  w\forces\varphi,\qquad
  \RR\subseteq\CC.
\]
The source of that line uses the meaningful names
\verb|\proves|, \verb|\entails|, \verb|\satisfies|,
\verb|\forces|, \verb|\RR|, and \verb|\CC|.
The full source-to-symbol-to-standard-LaTeX map is in
\texttt{docs/cheatsheet.md}.

Use \verb|\begin{DefinitionBlock}| to highlight a statement;
the unboxed \verb|definition| advances the same counter. A
\verb|ProofBlock| uses thin separator lines inside a theorem
rather than a second full frame. This lets long proofs cross
a page without being moved intact to the next sheet.
The abbreviations \WLoG, \WTS, \WRT, and \SFSoC{}
retain their small-cap styling.

\section{Prompt for transcribing handwritten notes}
Copy the text below into a chat when asking an LLM to typeset your
notes. The editable source is \path{docs/llm-prompt.txt}; the cheat
sheet points back to the same file.
\VerbatimInput[fontsize=\footnotesize]{docs/llm-prompt.txt}

\section{Citations and bibliography}
Footnotes are useful for comments on the text and can contain a
citation. For a short handout, keep references in \texttt{main.tex}
with \verb|\bibitem|. For an Angus-inspired
alphabetic bibliography, turn on the optional bibliography
setting in \texttt{andysnotes} and add a \texttt{references.bib}
file; see \path{docs/bibliography-biblatex.tex}. The same
\verb|\cite{key}| works in both cases. A journal or instructor's
required reference style takes precedence.

\clearpage
\appendix
\input{docs/symbols.tex}
\end{document}
